\section{ISRF channel: equations and discretisation}
\label{sec:isrf-eqdisc}

This section is a reference for the equations the GEAR interstellar-radiation-field
(ISRF) channel actually solves, and for the discretisation that solves them. It
is written to be read next to the source: every displayed equation is followed by
the file and line that computes it. It contains no validation status and no run
results.

\subsection{Scope, source tip, and where each piece runs}
\label{sec:isrf-eqdisc-scope}

The transported field is an M1 (first-moment, variable-Eddington-tensor) system,
two bands, three moments, one band-edge coupling between the two bands, and two
couplings out to the gas through Grackle. The scheme is \textbf{M1}, not the
``P1-relaxation'' some earlier comments named; that wording no longer occurs in
the source, and the file headers of
\code{radiation\_isrf.c} and \code{radiation\_propagation\_iact.h} say M1
(\code{radiation\_isrf.c:21-24}).

Line numbers below are from branch \code{darwin/gear\_radiation}. The inventory
this section was written from was read at tip \code{c122451718}; the citations
into \code{radiation\_isrf.c} past line~1114 have been re-anchored on the tip
\code{b07aee2fd}, which carries the relaxation-gate fix of
Section~\ref{sec:isrf-eqdisc-dissipation} (commit \code{6005e3c0f}). That commit
is the only change under \code{src/} between the two tips, so every other file's
line numbers are the same at both.

\paragraph{File keys.} Short basenames used in the citations, all relative to the
repository root:

\begin{center}
\begin{tabular}{p{4.6cm}p{9.9cm}}
\hline
\code{radiation\_isrf.c} & \code{src/feedback/GEAR/radiation\_isrf.c} \\
\code{radiation\_propagation\_iact.h} & \code{src/feedback/GEAR/radiation\_propagation\_iact.h} \\
\code{radiation\_iact.h} & \code{src/feedback/GEAR/radiation\_iact.h} \\
\code{radiation\_gas.c} & \code{src/feedback/GEAR/radiation\_gas.c} \\
\code{radiation.c} & \code{src/feedback/GEAR/radiation.c} \\
\code{radiation.h} & \code{src/feedback/GEAR/radiation.h} \\
\code{radiation\_table\_io.c} & \code{src/feedback/GEAR/radiation\_table\_io.c} \\
\code{stellar\_evolution.c} & \code{src/feedback/GEAR/stellar\_evolution.c} \\
\code{feedback\_struct.h} & \code{src/feedback/GEAR\_thermal/feedback\_struct.h} \\
\code{feedback\_properties.h} & \code{src/feedback/GEAR\_thermal/feedback\_properties.h} \\
\code{feedback.c} & \code{src/feedback/GEAR\_thermal/feedback.c} \\
\code{feedback\_iact.h} & \code{src/feedback/GEAR\_thermal/feedback\_iact.h} \\
\code{cooling\_gear\_subgrid.h} & \code{src/cooling/grackle/cooling\_gear\_subgrid.h} \\
\code{cooling\_properties.h}, \code{cooling.c}, \code{cooling\_io.h} & all in \code{src/cooling/grackle/} \\
\hline
\end{tabular}
\end{center}

\paragraph{Order within one SWIFT step, per gas particle.} The discretisation of
Section~\ref{sec:isrf-eqdisc-disc} is staggered across four task-graph stages, so
this ordering is needed to read it:

\begin{enumerate}
  \item \textbf{drift.} \code{feedback\_reset\_part} (\code{src/cell\_drift.c:357}
    into \code{feedback.c:226-233}) calls
    \code{radiation\_snapshot\_part\_propagation}
    (\code{radiation\_isrf.c:190}), which snapshots $u \to u_{\rm prev}$ and
    $\rho \to \rho_{\rm prev}$, sets $\kappa$, sets $c_{\rm hyp}$ for schemes
    0, 2 and 3, and draws the dose reservoir down; then the illumination-tag
    expiry (\code{radiation\_gas.c:557}); then
    \code{radiation\_cache\_m1\_closure\_part}
    (\code{radiation\_propagation\_iact.h:320}).
  \item \textbf{every density $h$-iteration.}
    \code{radiation\_init\_part\_propagation}
    (\code{radiation\_isrf.c:437-441}) zeroes the kernel-mean dissipation
    reference and resets the neighbour-bin maximum.
  \item \textbf{density loop.} \code{runner\_iact\_isrf\_propagation} and its
    non-symmetric twin (\code{radiation\_propagation\_iact.h:406},
    \code{:461}) accumulate the dissipation reference and
    \code{max\_ngb\_time\_bin}.
  \item \textbf{density ghost.} \code{radiation\_end\_density\_propagation}
    (\code{radiation\_isrf.c:493}) sets $c_{\rm hyp}$ for schemes 1 and 4 and
    rebuilds the closure cache (\code{radiation\_isrf.c:542}). The
    no-neighbour give-up path is \code{radiation\_part\_has\_no\_neighbours}
    (\code{radiation\_isrf.c:968}).
  \item \textbf{gradient loop.} \code{runner\_iact\_isrf\_gradient}
    (\code{radiation\_propagation\_iact.h:516}, \code{:570}) accumulates the
    M1 pressure-tensor divergence.
  \item \textbf{extra ghost.} \code{radiation\_end\_gradient\_propagation}
    (\code{radiation\_isrf.c:1232}) advances $\boldsymbol{F}$, applies the M1
    limiter, and sets both dissipation coefficients.
  \item \textbf{force loop.} \code{runner\_iact\_isrf\_dissipation}
    (\code{radiation\_propagation\_iact.h:634}, \code{:704}) accumulates
    $\nabla\!\cdot\!(\rho\boldsymbol{F}^{n+1})$ and the dissipation source.
  \item \textbf{end-force.} \code{radiation\_end\_force\_propagation}
    (\code{radiation\_isrf.c:768}, called from
    \code{src/runner\_others.c:863}) advances $u$.
\end{enumerate}

Star injection runs in the star feedback-apply loop,
\code{radiation\_iact\_nonsym\_feedback\_apply} (\code{radiation\_iact.h:202}).

\paragraph{The transport is off by default.}
\code{GEARFeedback:ISRF\_propagation} defaults to 0
(\code{feedback\_properties.h:1092-1093}), i.e.\ the default ISRF run is
injection plus receiver-side extinction only, with no transport. In that mode the
three ghosts return early (\code{radiation\_isrf.c:495}, \code{:770},
\code{:1235}) and both $\kappa$ are forced to zero
(\code{radiation\_isrf.c:216-218}), so nothing in
Sections~\ref{sec:isrf-eqdisc-continuum}, \ref{sec:isrf-eqdisc-disc} and
\ref{sec:isrf-eqdisc-transfer} is consumed; the pairwise density and gradient
loops nevertheless still run and still fill their accumulators
(\code{radiation\_isrf.c:206-211}).

\subsection{Symbols (reading aid)}
\label{sec:isrf-eqdisc-symbols}

A compact key, for reading this section only; the symbols are chosen here to
disambiguate letters the code reuses, and a paper would not need the table.

\begin{center}
\begin{tabular}{p{1.7cm}p{3.3cm}p{9.1cm}}
\hline
symbol & units & meaning \\
\hline
$u_m$ & energy/mass & specific band energy of moment $m$, \textbf{physical}, mass-specific, $a$-exponent 0 \\
$\boldsymbol{F}_m$ & energy/mass $\times$ velocity & specific band flux, \textbf{physical}; or its reduced form $\tilde{\boldsymbol{F}}_m = \boldsymbol{F}_m/c_{\rm hyp}$ under schemes 3 and 4 \\
$\boldsymbol{G}_m$ & energy/mass/length & the accumulator the code names \code{grad\_u}, physical, per physical length; see Section~\ref{sec:isrf-eqdisc-state} \\
$X_m$ & energy/mass/time & the accumulator \code{div\_specific\_flux}, i.e.\ $\rho^{-1}\nabla\!\cdot\!(\rho\boldsymbol{F}_m)$, physical \\
$\kappa_o$ & 1/length & local linear dust absorption rate of operator $o$, physical \\
$\kappa_{\rm eff}$ & area/mass & band dust \emph{mass} opacity, internal units \\
$c_{\rm hyp}$ & velocity & per-particle reduced propagation speed, physical \\
$c$ & velocity & true speed of light, \code{phys\_const->const\_speed\_light\_c} \\
$H$ & 1/time & Hubble rate, \code{cosmo->H} \\
$H_{\rm dil}$ & 1/time & dilated Hubble rate, $(c_{\rm hyp}/c)H$ \\
$\lambda_m$ & -- & band-edge weight of moment $m$ (Section~\ref{sec:isrf-eqdisc-lambda}) \\
$\mathcal{A}_m$ & -- & dimensionless relaxation depth of moment $m$ over one step \\
$\mathcal{E}_m$ & -- & one-step decay $e^{-\mathcal{A}_m}$ \\
$\mathcal{R}$ & -- & source rescale $c_{\rm hyp}/c$ \\
$\mathcal{C}_m$ & velocity$\,\times\,$time & flux-update coefficient, Eq.~(\ref{eq:isrfd-coeff-F}) \\
$\varphi$ & -- & integrating factor $(1-e^{-\mathcal{A}})/\mathcal{A}$ \\
$S_m$ & energy/mass/time & injected source rate, \code{u\_source\_rate} \\
$\Psi_m$ & energy/mass/time & artificial-dissipation source \\
$\mathsf{D}_o$ & -- & cached Eddington tensor of operator $o$ \\
$\hat f$ & -- & reduced flux $|\boldsymbol{F}|/(c_M u)$ feeding the closure \\
$a$ & -- & cosmological scale factor \\
$h$, $\mathcal{H}_i$ & length & smoothing length, \textbf{comoving} ($h_{\rm phys} = a h$); kernel support radius $\mathcal{H}_i = \gamma_K h_i$, comoving \\
$W'_i$ & length$^{-(d+1)}$ & kernel derivative, Eq.~(\ref{eq:isrfd-kernel-deriv}), comoving \\
$\omega_{ij}$ & mass & injection kernel weight, Eq.~(\ref{eq:isrfd-inj-weight}) \\
$\mathcal{X}$ & -- & band dust extinction factor, Eq.~(\ref{eq:isrfd-extinction}) \\
$\ell_{\rm ext}$ & length & extinction path, \textbf{comoving} \\
$\Sigma_{\rm gas}$ & mass/area & extinction column, \textbf{physical} \\
$\mathcal{T}$ & energy/mass & LW-to-PE band-edge transfer, Eq.~(\ref{eq:isrfd-transfer}) \\
$\mathcal{D}(Z)/\mathcal{D}_{\rm MW}$ & -- & dust-to-gas ratio relative to the Milky Way \\
$C_{\rm hyp}$ & -- & \code{GEARFeedback:ISRF\_c\_hyp\_margin}, default 0.5 \\
\hline
\end{tabular}
\end{center}

Throughout, ``physical'' and ``comoving'' are stated explicitly because the
pairwise operators take comoving inputs and write physical accumulators
(Section~\ref{sec:isrf-eqdisc-pairwise}).

\subsection{State and bands}
\label{sec:isrf-eqdisc-state}

\subsubsection{Moments and operators}

There are three \textbf{moments} (\code{feedback\_struct.h:51-56}):
\code{ISRF\_MOMENT\_PE} $=0$, \code{ISRF\_MOMENT\_LW} $=1$,
\code{ISRF\_MOMENT\_LW\_PHOTON} $=2$, with
\code{ISRF\_MOMENT\_COUNT} $=3$; and two \textbf{operators}
(\code{feedback\_struct.h:64-68}), \code{ISRF\_OPERATOR\_PE} $=0$ and
\code{ISRF\_OPERATOR\_LW} $=1$. The two maps are

\begin{align}
  o(m) &= \code{radiation\_isrf\_moment\_to\_operator}[m] = (\mathrm{PE},\,\mathrm{LW},\,\mathrm{LW}),
  \label{eq:isrfd-moment-to-op} \\
  m(o) &= \code{radiation\_isrf\_operator\_owner}[o] = (\mathrm{PE},\,\mathrm{LW}),
  \label{eq:isrfd-op-owner}
\end{align}

at \code{feedback\_struct.h:85-86} and \code{:115-116}, both declared unsized
with a \code{\_Static\_assert} on their length (\code{feedback\_struct.h:88-92},
\code{:118-122}). Equation~(\ref{eq:isrfd-moment-to-op}) means
\code{ISRF\_MOMENT\_LW\_PHOTON} shares LW's opacity, Eddington tensor and
dissipation coefficients; Equation~(\ref{eq:isrfd-op-owner}) means every write to
per-operator state reads its owning moment's $(u,\boldsymbol{F})$.

\code{ISRF\_MOMENT\_LW\_PHOTON} is a photon-\emph{number} moment carried in
energy-equivalent units: its source luminosity is a bare copy of the LW entry
(\code{stellar\_evolution.c:1401-1402} for the single-star path,
\code{:1674-1675} for the IMF-population path), and at first init the two are
forced equal whenever exactly one is zero, in either direction
(\code{radiation\_isrf.c:120-132}). It is \textbf{not} read by any coupling out
to the gas (Section~\ref{sec:isrf-eqdisc-couplings}); its only live role is the
diagnostic $d\ln(U/N)/dt$ identity, which needs its own band-edge weight
$\lambda_N(\mathrm{LW})$ (\code{radiation\_isrf.c:799-810}).

\subsubsection{Per-moment, per-operator and per-particle state}

Per moment, \code{struct feedback\_isrf\_moment\_data}
(\code{feedback\_struct.h:128-294}): $u$ (\textbf{double},
\code{feedback\_struct.h:165}), $u_{\rm prev}$ (\textbf{double}, \code{:173}),
\code{specific\_flux[3]} (float, \code{:193}), \code{div\_specific\_flux} (float,
\code{:204}), \code{dissipation\_u} (float, \code{:215}), \code{grad\_u[3]}
(float, \code{:222}), \code{u\_dose\_reservoir} (float, \code{:233}),
\code{u\_source\_rate} (float, \code{:243}), plus the debug-only
\code{u\_min\_since\_snapshot} (\code{:255}), \code{cumulative\_injected}
(\code{:268}) and \code{cumulative\_absorbed} (\code{:292}).

The double width of $u$ and $u_{\rm prev}$ is load-bearing and the code states why
(\code{feedback\_struct.h:157-164}): the per-step relaxation depth
$\mathcal{A}$ can be a few $10^{-8}$, and $1-\mathcal{A}$ at that scale is not
distinguishable in a 24-bit mantissa. Every other accumulator stays float because
none carries a multi-step decay at that depth.

\textbf{\code{grad\_u} is not the gradient of $u$.} The code accumulates the
anisotropic M1 pressure-tensor divergence
$\rho^{-1}\nabla\!\cdot\!(\mathsf{D}(\hat f)\,\rho\, u)$ into it
(\code{radiation\_propagation\_iact.h:373-388}, header comment \code{:334-336}).
The field name, and the ``$\nabla u$'' wording at
\code{radiation\_isrf.c:1210}, are a misnomer; the code is the truth. It is
written $\boldsymbol{G}_m$ here.

Under schemes 3 and 4, \code{specific\_flux} stores the \textbf{reduced} flux
$\tilde{\boldsymbol{F}} = \boldsymbol{F}/c_{\rm hyp}$, not the true flux
(\code{radiation\_isrf.c:1219-1227}, \code{radiation\_propagation\_iact.h:66}).

Per operator, \code{struct feedback\_isrf\_operator\_data}
(\code{feedback\_struct.h:302-361}): $\kappa$ (float, \code{:311}),
\code{m1\_closure\_D[3][3]} (float, \code{:319}), \code{ngb\_mean\_abs\_u\_V}
(float, \code{:326}), \code{dissipation\_alpha\_trigger} (float, \code{:342}) and
\code{dissipation\_alpha\_floor} (float, \code{:360}).

Per particle, in \code{struct feedback\_part\_data}: \code{max\_ngb\_time\_bin}
(\code{feedback\_struct.h:391}), the largest \code{time\_bin} seen in this
$h$-iteration's kernel; \code{rho\_prev} (float, \code{:442}), a
\textbf{comoving} density snapshot seeded to $1.0$ at first init and never left
at zero (\code{radiation\_isrf.c:158}, \code{:213}); \code{c\_hyp} (float,
\code{:478}), physical; \code{dt\_prev} (float, \code{:487}), this particle's own
\textbf{physical} timestep, cached raw and unfloored
(\code{radiation\_isrf.c:301}); and the four tag fields
\code{ISRF\_last\_touch\_ti}, \code{is\_illuminated\_ISRF},
\code{ISRF\_illumination\_end\_ti}, \code{ISRF\_reservoir\_end\_ti}
(\code{feedback\_struct.h:513}, \code{:532}, \code{:547}, \code{:556}).

On the star side, \code{struct feedback\_spart\_data.radiation} carries
\code{L\_band[ISRF\_MOMENT\_COUNT]} as \textbf{double} physical power in internal
units (\code{feedback\_struct.h:745}) and \code{Delta\_t} as float proper time
(\code{:759}).

\subsubsection{Band definitions and cross sections}
\label{sec:isrf-eqdisc-bands}

The photoelectric (PE) band is $6.0$--$11.2$\,eV and the Lyman-Werner (LW) band
is $11.2$--$13.6$\,eV. Only the lower edges enter the code as numbers:
$E_{\rm lo}(\mathrm{PE}) = 6.0$\,eV (\code{radiation.h:191}) and
$E_{\rm lo}(\mathrm{LW}) = 11.2$\,eV (\code{radiation.h:192}); the $13.6$\,eV
upper edge appears in the band comments (\code{radiation.h:104},
\code{feedback\_struct.h:41-43}) but is never used. The cgs forms are hard-coded
literals and never unit-converted (\code{radiation.h:202-203}):

\begin{equation}
  E_{\rm lo}(\mathrm{PE}) = 9.6130598\times10^{-12}\,{\rm erg}, \qquad
  E_{\rm lo}(\mathrm{LW}) = 1.7944378\times10^{-11}\,{\rm erg}.
  \label{eq:isrfd-band-edges}
\end{equation}

Dust cross sections per hydrogen nucleon (\code{radiation.h:105-106}), attributed
in the comment to \citet{Kim2023} and Weingartner \& Draine (2001)
(\code{radiation.h:102-104}):

\begin{equation}
  \sigma_d(\mathrm{PE}) = 9\times10^{-22}\,{\rm cm^2}, \qquad
  \sigma_d(\mathrm{LW}) = 1.5\times10^{-21}\,{\rm cm^2}.
  \label{eq:isrfd-sigma-d}
\end{equation}

The remaining constants are: $\mu_H = 1.4$, the mean mass per hydrogen nucleon in
proton masses (\code{radiation.h:100}), deliberately \emph{not} the run's own
hydrogen fraction; Grackle's solar metal fraction
$Z_\odot = 0.01295$ (\code{radiation.h:114}); the default dust-to-gas ratio
$f_{\rm gr} = 0.009387$ (\code{radiation.h:119}); and the Habing normalisation
$F_{\rm Habing} = 1.6\times10^{-3}\,{\rm erg\,s^{-1}\,cm^{-2}}$
(\code{radiation.h:123}).

For the H$_2$ photodissociation rate, the \textbf{primary} calibration constant is
the quotient

\begin{equation}
  \frac{\sigma_{\rm H_2}}{E_{\rm LW}} = 1.2847106348798106\times10^{-7}\,{\rm cm^2\,erg^{-1}},
  \label{eq:isrfd-sigma-over-e}
\end{equation}

at \code{radiation.h:147}, anchored on \citet{Sternberg2014} per its own comment
(\code{radiation.h:132-133}). The pair
$\sigma_{\rm H_2} = 2.5111667\times10^{-18}\,{\rm cm^2}$
(\code{radiation.h:167}) and $E_{\rm LW} = 12.2$\,eV (\code{radiation.h:154}) is
derived, quoted only for provenance, and is \emph{not} read by the rate.

\subsubsection{Band-edge weights $\lambda_E(\mathrm{PE})$, $\lambda_E(\mathrm{LW})$, $\lambda_N(\mathrm{LW})$}
\label{sec:isrf-eqdisc-lambda}

Each $\lambda_m$ multiplies the dilated Hubble rate inside its own moment's
relaxation depth (\code{radiation\_isrf.c:897}, \code{:1289}). Its meaning, as
the code uses it: the number of e-folds-worth of band content lost per e-fold of
expansion, because the band edges are fixed in \emph{physical} energy while
photons redshift downward through them.

The grey (comoving-edge) value \textbf{differs per moment}:

\begin{align}
  \lambda_E(b) &= 1 + \frac{\Lambda_b\, E_{\rm lo}(b)}{\langle E\rangle_b},
  \qquad b \in \{\mathrm{PE}, \mathrm{LW}\},
  \label{eq:isrfd-lambda-E} \\
  \lambda_N(\mathrm{LW}) &= \Lambda_{\rm LW}
  \qquad \text{(no $+1$)},
  \label{eq:isrfd-lambda-N}
\end{align}

at \code{radiation.h:225} and \code{radiation.c:91}, \code{:222-224}. The $+1$ in
Equation~(\ref{eq:isrfd-lambda-E}) is the redshift loss of photon \emph{energy},
which survives when the edges are comoving, so the grey value of $\lambda_E$ is
$1$. Photon \emph{number} does not redshift, so the grey value of
$\lambda_N(\mathrm{LW})$ is $0$, not $1$ (rationale
\code{radiation.c:226-230}). The three shipped fallbacks are mutually consistent
on this: $(6.508 - 1)\times 12.2/11.2 = 6.0$ (\code{radiation.h:237-239}).

The three scalars are computed once at start-up, over the whole IMF
$[\code{imf.mass\_min}, \code{imf.mass\_max}]$, at one reference metallicity
$Z_{\rm ref} = 0.014$ for a 2D table (\code{radiation.h:176},
\code{radiation.c:122-127}):

\begin{align}
  \lambda_E(\mathrm{PE}) &= 1 + \frac{l_{\rm edge}(\mathrm{PE})}{l_{\rm PE}},
  \label{eq:isrfd-lambda-pe-table} \\
  \lambda_E(\mathrm{LW}) &= 1 + \frac{l_{\rm edge}(\mathrm{LW})}{l_{\rm LW}},
  \label{eq:isrfd-lambda-lw-table} \\
  \lambda_N(\mathrm{LW}) &= \big(\lambda_E(\mathrm{LW}) - 1\big)
    \frac{\langle E\rangle_{\rm LW}}{E_{\rm lo}(\mathrm{LW})},
  \label{eq:isrfd-lambda-n-table}
\end{align}

at \code{radiation.c:212}, \code{:213} and \code{:221-224}, with the ratio itself
formed at \code{radiation.c:165-166}. Here $l_{\rm edge}(b)$ is the
IMF-integrated \code{Integrated\_SpectralPhotonRateAt<b>Edge} differenced over
the mass interval, \emph{already} multiplied by $E_{\rm lo}(b)^2$ and converted to
internal power at read time (the read divides the cgs value by
$\code{units\_cgs(POWER)}/E_{\rm lo}(b)^2$,
\code{radiation\_table\_io.c:1232-1234} for PE and \code{:1264-1266} for LW; the
division itself is \code{radiation\_table\_io.c:668}), so both sides of the ratio
are internal power and no further factor is applied
(\code{radiation.c:161-165}). $l_b$ is \code{Integrated\_L\_PE} or
\code{Integrated\_L\_LW} over the same mass interval, and
$\langle E\rangle_{\rm LW}$ is \code{Integrated\_MeanPhotonEnergyLW} read at
\code{mass\_max} as a single point, not differenced
(\code{radiation.c:147-151}), in cgs erg.

Guards and tripwires: with radiation inactive or the ISRF off, the three keep
their compile-time fallbacks (\code{radiation.c:115-120}); so do they if
$l_{\rm PE} \le 0$ or $l_{\rm LW} \le 0$, tested exactly with no epsilon
(\code{radiation.c:159}). A zero numerator over a nonzero denominator gives
exactly grey, $\lambda = 1$, with a rank-0 warning
(\code{radiation.c:176-189}). A positive but implausibly small
$\lambda - 1 < 0.05$ is fatal (\code{radiation.c:197-210}, macro
\code{radiation.c:60}); the same floor is applied to $\lambda_N(\mathrm{LW})$
\emph{itself}, not to $\lambda_N - 1$ (\code{radiation.c:231-239}, rationale
\code{:226-230}). A value above 250 is a warning only
(\code{radiation.c:61}). The compile-time fallbacks, described as the
young-population end of a one-parameter spectral family, are
$\lambda_E(\mathrm{PE}) = 2.154$, $\lambda_E(\mathrm{LW}) = 6.508$,
$\lambda_N(\mathrm{LW}) = 6.0$ (\code{radiation.h:237-239}), stored as double
fields of \code{feedback\_props} (\code{feedback\_properties.h:252-254}).

\subsection{The continuum M1 system, per band}
\label{sec:isrf-eqdisc-continuum}

The code never writes the PDE; the system below is the continuum problem whose
one-step solution is implemented in Sections~\ref{sec:isrf-eqdisc-u-update} and
\ref{sec:isrf-eqdisc-F-update}. Per moment $m$, with operator $o = o(m)$:

\begin{align}
  \frac{\partial u_m}{\partial t} &=
    -\big(c_{\rm hyp}\kappa_o + \lambda_m H_{\rm dil}\big)\, u_m
    + \frac{c_{\rm hyp}}{c}\, S_m
    - \frac{1}{\rho}\nabla\!\cdot\!\big(\rho \boldsymbol{F}_m\big)
    + \Psi_m,
  \label{eq:isrfd-continuum-u} \\
  \frac{\partial \boldsymbol{F}_m}{\partial t} &=
    -\big(c_{\rm hyp}\kappa_o + \lambda_m H_{\rm dil}\big)\, \boldsymbol{F}_m
    - c_{\rm hyp}^2\, \frac{1}{\rho}\nabla\!\cdot\!
      \big(\mathsf{D}_o(\hat f)\, \rho\, u_m\big),
  \label{eq:isrfd-continuum-F}
\end{align}

with

\begin{equation}
  H_{\rm dil} = \frac{c_{\rm hyp}}{c}\, H .
  \label{eq:isrfd-Hdil}
\end{equation}

at \code{radiation\_isrf.c:789} (double) and \code{:1249} (float).

Equations~(\ref{eq:isrfd-continuum-u}) and (\ref{eq:isrfd-continuum-F}) are
not written anywhere in the source; they are implemented by the one-step updates
of Equations~(\ref{eq:isrfd-u-update}) and (\ref{eq:isrfd-F-update}), at
\code{radiation\_isrf.c:933-937} and \code{:1308}, with the depths of
Equations~(\ref{eq:isrfd-depth-u}) and (\ref{eq:isrfd-depth-F}) at
\code{radiation\_isrf.c:897} and \code{:1289}.

The one-step solution is \emph{exact} with $S_m$ and
$\nabla\!\cdot\!(\rho\boldsymbol{F}_m)$ frozen, but \textbf{not} including the
artificial dissipation $\Psi_m$, which is decayed together with $u_{\rm prev}$
rather than integrated as a frozen source (\code{radiation\_isrf.c:682-683},
placement rationale \code{:697-700}).

\paragraph{Closure.} \code{radiation\_build\_m1\_closure\_tensor}
(\code{radiation\_propagation\_iact.h:275-284}), with coefficients from
\code{radiation\_get\_m1\_closure\_coefficients\_band}
(\code{radiation\_propagation\_iact.h:247-262}):

\begin{align}
  \hat f &= \min\!\left(1, \frac{|\boldsymbol{F}|}{c_M u}\right)
    \ \ \text{for } c_M u > 0, \quad \hat f = 0 \text{ otherwise},
  \label{eq:isrfd-fhat} \\
  \chi &= \frac{3 + 4\hat f^{\,2}}{5 + 2\sqrt{4 - 3\hat f^{\,2}}},
  \label{eq:isrfd-chi} \\
  \mathsf{D} &= \frac{1-\chi}{2}\,\mathsf{I}
    + \frac{3\chi - 1}{2}\, \boldsymbol{n}\otimes\boldsymbol{n},
    \qquad \boldsymbol{n} = \frac{\boldsymbol{F}}{|\boldsymbol{F}|},
    \ \ \boldsymbol{n} = 0 \text{ at } \boldsymbol{F} = 0 .
  \label{eq:isrfd-eddington}
\end{align}

Here $c_M = c_{\rm hyp}$, or exactly $1$ under schemes 3 and 4
(\code{radiation\_propagation\_iact.h:323},
\code{radiation\_isrf.c:1243}), because the stored flux is already reduced.

\paragraph{Absorption.} Pure dust absorption: no scattering, and no emission term
other than the star source. $\kappa_o$ is the band mass opacity times the local
physical density (\code{radiation\_isrf.c:1577-1579}), with no column.
Section~\ref{sec:isrf-eqdisc-extinction}'s extinction is a separate,
independently computed quantity and the code states the two are not
interchangeable (\code{feedback\_struct.h:308-310}).

\paragraph{Cosmological terms.} Exactly \emph{one} power of $H$ appears, in both
equations. The code's justification (\code{radiation\_isrf.c:702-705}): $u$ is
\emph{mass}-specific, so the volume dilution is already carried by the physical
gas density it is measured against, and the specific flux loses the same single
power. $H$ is \textbf{dilated} by the same $c_{\rm hyp}/c$ factor as the
absorption term, in double at \code{radiation\_isrf.c:789} and in float at
\code{radiation\_isrf.c:1249}. The reason given
(\code{radiation\_isrf.c:681-691}): unlike $\kappa$, $H$ does not pick that
factor up for free, and leaving it bare suppresses the true-speed fixed point of
the homogeneous equation and stops $c_{\rm hyp}$ cancelling out of it. A
documented consequence, not a defect (\code{radiation\_isrf.c:712-717}): with
$\kappa = 0$ the pure-expansion transient decays at $(c_{\rm hyp}/c)H$, not at
$H$. The term vanishes exactly in a non-cosmological run, because SWIFT sets
\code{cosmo->H} $=0$.

The Hubble term is folded \emph{into} the relaxation depth rather than subtracted
explicitly, so $e^{-\mathcal{A}}$ cannot drive $u$ negative at large $H\,dt$
(\code{radiation\_isrf.c:706-711}).

\paragraph{Comoving bookkeeping of the operators.} Inputs are comoving,
accumulators are physical, and each spatial operator carries one net inverse
length, so each closes with \code{a\_factor\_comoving\_to\_physical} $=1/a$
(\code{radiation\_propagation\_iact.h:38-45}, applied at \code{:541},
\code{:597}, \code{:664}, \code{:734}).

\subsection{The discretisation}
\label{sec:isrf-eqdisc-disc}

\subsubsection{Pairwise transport operators}
\label{sec:isrf-eqdisc-pairwise}

Throughout, $\boldsymbol{dx} = \boldsymbol{x}_i - \boldsymbol{x}_j$ and
$r = |\boldsymbol{dx}|$ are \textbf{comoving}, $r^{-1} = 1/r$, every $\rho$ is the
\textbf{comoving} \code{rho\_prev} snapshot, and the kernel derivative the code
calls \code{wi\_dr} is written
\begin{equation}
  W'_i \equiv \frac{dW}{dq}\!\left(\frac{r}{h_i}\right) h_i^{-(d+1)},
  \qquad
  W'_j \equiv \frac{dW}{dq}\!\left(\frac{r}{h_j}\right) h_j^{-(d+1)},
  \label{eq:isrfd-kernel-deriv}
\end{equation}
with $h$ comoving and $d$ the dimension count.

\paragraph{Divergence (force loop).}
\code{radiation\_divergence\_accumulate\_band}
(\code{radiation\_propagation\_iact.h:91-109}):

\begin{align}
  \Phi_{ij} &= \left(
    \frac{\boldsymbol{F}_i\!\cdot\!\boldsymbol{dx}}{\rho_i}\,W'_i\,r^{-1}
    + \frac{\boldsymbol{F}_j\!\cdot\!\boldsymbol{dx}}{\rho_j}\,W'_j\,r^{-1}
  \right)\frac{1}{a},
  \label{eq:isrfd-div-phi} \\
  X_i &\mathrel{+}= + m_j \Phi_{ij}, \qquad
  X_j \mathrel{+}= - m_i \Phi_{ij},
  \label{eq:isrfd-div-apply}
\end{align}

at \code{radiation\_propagation\_iact.h:94-96} and \code{:108-109}, where $X$
denotes \code{div\_specific\_flux}. Under schemes 3 and 4 the applied form is
instead $X_i \mathrel{+}= + c_i m_j \Phi_{ij}$ and
$X_j \mathrel{+}= - c_j m_i \Phi_{ij}$
(\code{radiation\_propagation\_iact.h:103-104}), with $c_i$ the receiving
particle's own $c_{\rm hyp}$.

\textbf{What is conserved exactly.} At the \emph{accumulator} level, per pair,
$m_i X_i + m_j X_j = 0$ identically, because both sides share one $\Phi_{ij}$
(\code{radiation\_propagation\_iact.h:108-109}). Under schemes 3 and 4 the
conserved combination is instead $\sum_i m_i X_i / c_{{\rm hyp},i}$, which the
comment states explicitly (\code{radiation\_propagation\_iact.h:100-102},
\code{:206-207}). The exactness is a property of the accumulator, not of the
applied update: $u$ is advanced by $dt_i\,\varphi_i\,X_i$ per particle
(\code{radiation\_isrf.c:933-937}), so the pairwise cancellation survives into
$u$ only where $dt_i = dt_j$ and $\varphi_i = \varphi_j$, i.e.\ same time bin and
same relaxation depth. Across time bins it does not.

The divergence sits in the \emph{force} loop, not a type-1 loop, because it needs
both sides of every pair credited and only the type-2 dispatch fires both sides
whenever either kernel reaches (\code{radiation\_propagation\_iact.h:26-33},
\code{feedback\_struct.h:195-199}).

\paragraph{Gradient (gradient loop).}
\code{radiation\_gradient\_accumulate\_band}
(\code{radiation\_propagation\_iact.h:368-389}):

\begin{align}
  \boldsymbol{T}_i &= \big(\mathsf{D}_i\!\cdot\!\boldsymbol{dx}\big)\,\rho_i u_i\, r^{-1},
  \qquad
  \boldsymbol{T}_j = \big(\mathsf{D}_j\!\cdot\!\boldsymbol{dx}\big)\,\rho_j u_j\, r^{-1},
  \label{eq:isrfd-grad-temp} \\
  \mathcal{F}_i &= \frac{m_j}{\rho_i^2}\,W'_i\,\frac{1}{a},
  \qquad
  \mathcal{F}_j = \frac{m_i}{\rho_j^2}\,W'_j\,\frac{1}{a},
  \label{eq:isrfd-grad-fac} \\
  \boldsymbol{G}_i &\mathrel{+}= -\big(\boldsymbol{T}_i - \boldsymbol{T}_j\big)\mathcal{F}_i,
  \qquad
  \boldsymbol{G}_j \mathrel{+}= -\big(\boldsymbol{T}_i - \boldsymbol{T}_j\big)\mathcal{F}_j,
  \label{eq:isrfd-grad-apply}
\end{align}

at \code{radiation\_propagation\_iact.h:377-378}, \code{:381-384} and
\code{:387-388}. Each particle uses its \emph{own} kernel derivative and its
\emph{own} closure tensor, with no shared average and no grad-$h$ \code{forcef}
factor. The comment states that this is the deliberate complement of the
divergence's shared-coefficient form, and that pairing the two makes them exactly
skew-adjoint in the $m\rho$ inner product
(\code{radiation\_propagation\_iact.h:338-341},
\code{feedback\_struct.h:429-435}); that algebra is reported here as the code's
own claim and has not been re-derived.

Both operators read \code{rho\_prev}, never the live \code{p->rho}
(\code{radiation\_propagation\_iact.h:42-45}, \code{:423-424}, \code{:536-537},
\code{:654-655}): during the density loop \code{p->rho} is a partial sum, and the
skew-adjoint pairing needs both operators built from the same densities.

\paragraph{Dissipation reference (density loop).}
\code{radiation\_dissipation\_reference\_accumulate\_band}
(\code{radiation\_propagation\_iact.h:142-143}):

\begin{align}
  \code{ngb\_mean\_abs\_u\_V}_i &\mathrel{+}= \frac{m_j}{\rho_j}\, w_i\,
    \big|\rho_j\, u_{{\rm prev},j}\big|,
  \label{eq:isrfd-diss-ref-i} \\
  \code{ngb\_mean\_abs\_u\_V}_j &\mathrel{+}= \frac{m_i}{\rho_i}\, w_j\,
    \big|\rho_i\, u_{{\rm prev},i}\big| .
  \label{eq:isrfd-diss-ref-j}
\end{align}

There is deliberately no $1/a$ factor: its only consumer divides it into
$\rho_{\rm prev} u$, which carries the same $a^3$ weight
(\code{radiation\_propagation\_iact.h:120-123}).

\paragraph{Dissipation source (force loop).}
\code{radiation\_dissipation\_force\_accumulate\_band}
(\code{radiation\_propagation\_iact.h:191-216}):

\begin{align}
  d_{ij} &= \rho_i u_i - \rho_j u_j,
  \label{eq:isrfd-diss-d} \\
  \alpha_{ij} &= \max\big(\alpha^{\rm trig}_i,\ \alpha^{\rm trig}_j,\
    \alpha^{\rm floor}_i,\ \alpha^{\rm floor}_j\big),
  \label{eq:isrfd-diss-alpha} \\
  \overline{W}_{ij} &= \tfrac{1}{2}\big(W'_i + W'_j\big),
  \label{eq:isrfd-diss-wbar} \\
  \Upsilon_{ij} &= \frac{d_{ij}\,\overline{W}_{ij}}{\rho_i \rho_j}\,\frac{1}{a},
  \label{eq:isrfd-diss-shape} \\
  \Psi_{ij} &= \alpha_{ij}\,\min(c_i, c_j)\,\Upsilon_{ij},
  \label{eq:isrfd-diss-psi} \\
  \code{dissipation\_u}_i &\mathrel{+}= + m_j \Psi_{ij},
  \qquad
  \code{dissipation\_u}_j \mathrel{+}= - m_i \Psi_{ij},
  \label{eq:isrfd-diss-apply}
\end{align}

at \code{radiation\_propagation\_iact.h:191}, \code{:196-198}, \code{:200},
\code{:202-203}, \code{:213} and \code{:215-216}. Under schemes 3 and 4 the two
applied terms are $+ m_j \alpha_{ij} c_i \Upsilon_{ij}$ and
$- m_i \alpha_{ij} c_j \Upsilon_{ij}$
(\code{radiation\_propagation\_iact.h:208-209}), i.e.\ two receiver-side speeds
rather than one shared minimum. The accumulator-level antisymmetry, and the
caveat about the applied update, are the same as for the divergence.

The dissipation is applied to a pair \emph{ungated}, with no mutual-reach test,
because the force loop fires both sides whenever either kernel reaches
(\code{radiation\_propagation\_iact.h:154-155}). All four operators disable
floating-point reassociation under clang for the shared blocks
(\code{radiation\_propagation\_iact.h:88-90}, \code{:188-190}) so the two scheme
families stay bit-identical at uniform $c_{\rm hyp}$; the comment notes GCC has no
block-scoped equivalent.

\subsubsection{The exact-exponential relaxation update of $u$ (end-force ghost)}
\label{sec:isrf-eqdisc-u-update}

\code{radiation\_end\_force\_propagation},
\code{radiation\_isrf.c:768-951}. With $dt = \code{dt\_prev}$:

\begin{align}
  \mathcal{R} &= \frac{c_{\rm hyp}}{c},
  \qquad H_{\rm dil} = \mathcal{R}\, H,
  \label{eq:isrfd-rescale} \\
  \mathcal{A}_m &= \big(c_{\rm hyp}\kappa_o + \lambda_m H_{\rm dil}\big)\, dt,
  \label{eq:isrfd-depth-u} \\
  \mathcal{E}_m &= e^{-\mathcal{A}_m},
  \qquad
  \varphi_m = \frac{1 - e^{-\mathcal{A}_m}}{\mathcal{A}_m},
  \label{eq:isrfd-decay-u}
\end{align}

at \code{radiation\_isrf.c:775}, \code{:789}, \code{:897}, \code{:898} and
\code{:899}. The update itself is

\begin{equation}
  u_m^{n+1} = \mathcal{E}_m\Big(u_{{\rm prev},m}
    + dt\,\varphi_m\,\code{dissipation\_u}_m\Big)
    + dt\,\varphi_m\Big(\mathcal{R}\,S_m - X_m\Big),
  \label{eq:isrfd-u-update}
\end{equation}

at \code{radiation\_isrf.c:933-937}, plus, for $m = \code{ISRF\_MOMENT\_PE}$
only, a raw add of the band-edge transfer of
Section~\ref{sec:isrf-eqdisc-transfer},

\begin{equation}
  u_{\rm PE}^{n+1} \mathrel{+}= \mathcal{T},
  \label{eq:isrfd-u-transfer-add}
\end{equation}

written as a separate statement at \code{radiation\_isrf.c:944} so that
\code{-ffast-math} cannot merge it into a form that no longer vanishes exactly at
$H = 0$ (\code{radiation\_isrf.c:939-943}).

Both $dt$ and $c_{\rm hyp}$ are read from the particle and never recomputed: the
call site's own $dt$ uses a different step-begin convention, and using it would
desynchronise this update from the flux update, so the function takes no $dt$
argument at all (\code{radiation\_isrf.c:758-763}).

The integrating factor is \code{radiation\_relaxation\_phi\_factor}
(float, \code{radiation\_isrf.c:562-565}) and
\code{radiation\_relaxation\_phi\_factor\_double}
(\code{radiation\_isrf.c:580-583}):

\begin{equation}
  \varphi(\mathcal{A}) =
  \begin{cases}
    1 - \dfrac{\mathcal{A}}{2} + \dfrac{\mathcal{A}^2}{6}, & \mathcal{A} < 10^{-6},\\[2ex]
    -\dfrac{\code{expm1}(-\mathcal{A})}{\mathcal{A}}, & \text{otherwise},
  \end{cases}
  \label{eq:isrfd-phi}
\end{equation}

at \code{radiation\_isrf.c:563-564} and \code{:581-582}. As
$\mathcal{A}\to 0$, $\varphi \to 1$ and $\mathcal{E}\to 1$, so
Equation~(\ref{eq:isrfd-u-update}) degenerates to plain forward Euler transport,
$u^{n+1} = u_{\rm prev} + dt(S - X + \Psi)$. As $\mathcal{A}\to\infty$,
$\varphi \to 1/\mathcal{A}$, i.e.\ exact relaxation to the local target each
step. The $10^{-6}$ Taylor branch exists to avoid the $0/0$ cancellation
(\code{radiation\_isrf.c:550-553}).

$u$ is advanced in \textbf{double} throughout
(\code{radiation\_isrf.c:773-789}, \code{:897-899}, \code{:933-937}); the double
$\varphi$ exists specifically because the $u$-path depth can be a few $10^{-8}$
(\code{radiation\_isrf.c:569-579}). The update is \emph{idempotent}: $u$ is
rebuilt from $u_{\rm prev}$ plus this step's accumulators, never incremented
(\code{radiation\_isrf.c:745-748}), the two debug ledger counters being the
stated exception (\code{radiation\_isrf.c:916-930}). No flux limiter is applied
here: the extra ghost already limited this step's flux against $u^n$
(\code{radiation\_isrf.c:746-748}).

\subsubsection{The exact-exponential relaxation update of $\boldsymbol{F}$ (extra ghost)}
\label{sec:isrf-eqdisc-F-update}

\code{radiation\_end\_gradient\_propagation},
\code{radiation\_isrf.c:1232-1380}. Everything here is \textbf{float},
including $H_{\rm dil}$ and the \code{expf}:

\begin{align}
  c_M &= 1 \ \text{under schemes 3/4, else } c_{\rm hyp},
  \label{eq:isrfd-cM} \\
  H_{\rm dil} &= \frac{c_{\rm hyp}}{c}\, H,
  \label{eq:isrfd-Hdil-float} \\
  \mathcal{A}_m &= \big(c_{\rm hyp}\kappa_o + \lambda_m H_{\rm dil}\big)\, dt,
  \qquad
  \mathcal{E}_m = e^{-\mathcal{A}_m},
  \qquad
  \varphi_m = \varphi(\mathcal{A}_m),
  \label{eq:isrfd-depth-F} \\
  \mathcal{C}_m &=
  \begin{cases}
    c_{\rm hyp}\, dt\, \varphi_m, & \text{schemes 3, 4},\\
    c_{\rm hyp}^2\, dt\, \varphi_m, & \text{schemes 0, 1, 2},
  \end{cases}
  \label{eq:isrfd-coeff-F} \\
  \boldsymbol{F}_m^{n+1} &= \mathcal{E}_m \boldsymbol{F}_m
    - \mathcal{C}_m\, \boldsymbol{G}_m,
  \label{eq:isrfd-F-update}
\end{align}

at \code{radiation\_isrf.c:1243}, \code{:1249}, \code{:1289-1291},
\code{:1294-1296} and \code{:1308}.

\textbf{Precision asymmetry}, recorded as a property of the scheme. The same
moment's $u$ and $\boldsymbol{F}$ relax with depths computed at different
precision, double at \code{radiation\_isrf.c:897} and float at
\code{radiation\_isrf.c:1289}, from the same inputs. At the few-$10^{-8}$ depths
the $u$-path doxygen describes, the float depth is unresolvable, so the flux
decay is effectively 1 there while the energy decay is not. The code's own
argument for keeping the flux path float
(\code{radiation\_isrf.c:569-575}) is that $\boldsymbol{F}$ itself is float, so a
double $\varphi$ there would not change what gets stored. That sits in tension
with \code{radiation\_isrf.c:1262-1267}, which requires $\boldsymbol{F}$ and $u$
of one moment to relax at the \emph{same} rate, or the reduced flux
$\hat f = |\boldsymbol{F}|/(c_{\rm hyp}u)$ that the closure reads is corrupted.

\paragraph{M1 flux limiter.} Decided once per \emph{operator} from its owning
moment (\code{radiation\_isrf.c:1327-1336}) and applied identically to every
moment sharing that operator (\code{radiation\_isrf.c:1377-1378}):

\begin{equation}
  \boldsymbol{F} \longleftarrow
  \begin{cases}
    0, & u \le 0,\\
    \boldsymbol{F} \ \text{(untouched, no multiply)}, & |\boldsymbol{F}|^2 \le 0,\\
    \boldsymbol{F}\,\min\!\left(1, \dfrac{c_M u}{|\boldsymbol{F}|}\right), & \text{otherwise},
  \end{cases}
  \label{eq:isrfd-limiter}
\end{equation}

at \code{radiation\_isrf.c:628}, \code{:632}, \code{:634} for the decision and
\code{:656-667} for the application.
$\boldsymbol{F}\!\cdot\!\boldsymbol{F}$ and the ratio are formed in
\textbf{double}, because in float32
$\boldsymbol{F}\!\cdot\!\boldsymbol{F}$ underflows once
$|\boldsymbol{F}| \lesssim \sqrt{\code{FLT\_MIN}} \approx 1.1\times10^{-19}$
internal, which would skip the limiter for a nonzero flux
(\code{radiation\_isrf.c:630-634}, rationale \code{:610-613}). The three outcomes
are a \code{switch}, not one scale value, so the skip case executes no multiply at
all under \code{-ffast-math} (\code{radiation\_isrf.c:655-668}, rationale
\code{:641-645}).

\code{div\_specific\_flux} is zeroed here, not at the drift, so a snapshot keeps
the last step's value (\code{radiation\_isrf.c:1314}).

\subsubsection{Closure cache}

\code{radiation\_cache\_m1\_closure\_part}
(\code{radiation\_propagation\_iact.h:320-331}) stores one $\mathsf{D}$ per
operator, built from that operator's \emph{owning} moment's
$(u, \code{specific\_flux})$ and its $c_M$. $|\boldsymbol{F}|^2$,
$|\boldsymbol{F}|$, $c_M u$ and $\hat f$ are formed in double for the same
underflow reason (\code{radiation\_propagation\_iact.h:247-256}), and zero guards
keep $\boldsymbol{F} = 0$ at $n=0$ giving $\hat f = 0$ rather than a NaN
(\code{radiation\_propagation\_iact.h:249}). Three call sites: first init
(\code{radiation\_isrf.c:163}), the drift reset (\code{feedback.c:232}, after the
tag expiry, which can zero $u$), and the density ghost once this step's
$c_{\rm hyp}$ is known (\code{radiation\_isrf.c:542}).

\subsubsection{Artificial dissipation: trigger, floor, gate}
\label{sec:isrf-eqdisc-dissipation}

Defaults, all parsed unconditionally so a validation run can set them with
propagation off: \code{ISRF\_dissipation\_alpha\_max} $=0.5$
(\code{feedback\_properties.h:1199-1200}),
\code{ISRF\_dissipation\_negativity\_threshold} $\epsilon_1 = 0.01$
(\code{:1201-1202}), \code{ISRF\_dissipation\_alpha\_floor} $=0.5$
(\code{:1221-1222}), \code{ISRF\_dissipation\_floor\_h\_over\_lambda}
$\epsilon_\lambda = 0.5$ (\code{:1223-1224}),
\code{ISRF\_dissipation\_floor\_relaxation\_residual} $\epsilon_R = 0.40$
(\code{:1240-1242}) and \code{ISRF\_dissipation\_alpha\_pin\_for\_debugging}
$=0$ (\code{:1203-1204}). Elsewhere: \code{ISRF\_c\_hyp\_margin} $=0.5$
(\code{:1169-1170}), with \code{ISRF\_c\_hyp\_fixed\_fraction\_of\_c} and
\code{ISRF\_c\_hyp\_pin\_for\_debugging} both disabled by default. Ranges are
enforced at parse time: $\alpha_{\max} \ge 0$ (\code{:1206-1208}),
$\epsilon_1 \in (0,1]$ (\code{:1210-1215}), $\alpha_{\rm floor} \ge 0$
(\code{:1226-1230}), $\epsilon_\lambda > 0$ (\code{:1232-1236}),
$\epsilon_R \in [0,1]$ (\code{:1244-1252}).

\paragraph{Trigger.} \code{radiation\_update\_dissipation\_alpha\_band}
(\code{radiation\_isrf.c:1005-1028}), evaluated once per operator per step in the
extra ghost (\code{radiation\_isrf.c:1355}) with
$u_V = \rho_{\rm prev}\, u$ (\code{radiation\_isrf.c:1338}):

\begin{align}
  u_V^{\rm ref} &= \max\big(\code{ngb\_mean\_abs\_u\_V},\ -u_V\big),
  \label{eq:isrfd-trig-ref} \\
  \epsilon &=
    \begin{cases}
      -u_V / u_V^{\rm ref}, & u_V < 0 \ \text{and}\ u_V^{\rm ref} > 0,\\
      0, & \text{otherwise},
    \end{cases}
  \label{eq:isrfd-trig-eps} \\
  x &= \min\!\left(\frac{\epsilon}{\epsilon_1},\, 1\right),
  \qquad
  \alpha^{\rm aim} = \alpha_{\max}\, x^2 (3 - 2x),
  \label{eq:isrfd-trig-aim} \\
  \alpha^{\rm trig} &=
    \begin{cases}
      \alpha^{\rm aim}, & \alpha^{\rm aim} \ge \alpha^{\rm trig}_{\rm prev},\\
      \alpha^{\rm aim} + \big(\alpha^{\rm trig}_{\rm prev} - \alpha^{\rm aim}\big)\mathcal{E}_\alpha,
        & \text{otherwise},
    \end{cases}
  \label{eq:isrfd-trig-apply} \\
  \mathcal{E}_\alpha &= \exp\!\left(
    -\frac{c_{\rm hyp}\, dt}{5.0\, h_{\rm phys}} - c_{\rm hyp}\kappa\, dt\right),
  \label{eq:isrfd-trig-decay}
\end{align}

at \code{radiation\_isrf.c:1016}, \code{:1017}, \code{:1018-1019},
\code{:1021}, \code{:1027} and \code{:1023-1026}, with
$5.0 = \code{RADIATION\_ISRF\_DISSIPATION\_DECAY\_LENGTH}$
(\code{radiation.h:70}). Instant rise, exponential decay. The decay memory is the
trigger's \emph{own} previous value, not the combined coefficient, so a high floor
cannot hold up the trigger's decay tail
(\code{radiation\_isrf.c:1351-1356}). The trigger responds to
\textbf{negativity only} and is exactly zero on a positive front
(\code{radiation\_isrf.c:1032-1036}). The test
$u_V < 0 \wedge u_V^{\rm ref} > 0$ in
Equation~(\ref{eq:isrfd-trig-eps}) is exact rather than assumed, because
$\rho_{\rm prev}u$ can narrow to $-0$ for a faint band while the unnarrowed
double still reads negative, which would make the division $0/0$
(\code{radiation\_isrf.c:1010-1015}).

\paragraph{Floor.} \code{radiation\_dissipation\_alpha\_floor\_band}
(\code{radiation\_isrf.c:1061-1063}):

\begin{equation}
  x = \frac{h_{\rm phys}\,\kappa}{\epsilon_\lambda},
  \qquad
  \alpha^{\rm floor} = \frac{\alpha_{\rm floor}}{1 + x^4}.
  \label{eq:isrfd-floor}
\end{equation}

The roll-off is quartic, not quadratic, so an optically thick region
($h\kappa \sim 6$--$10$) gets about $10^{-5}$ of the floor rather than a
percent-level residue (\code{radiation\_isrf.c:1040-1045}). The floor enters a
pair unconditionally through Equation~(\ref{eq:isrfd-diss-alpha}); this roll-off
and the gate below are its only gating (\code{radiation\_isrf.c:1047-1048}).

\paragraph{Relaxation-residual gate.}
\code{radiation\_dissipation\_floor\_relaxation\_gate}
(\code{radiation\_isrf.c:1128-1187}), which multiplies the floor's aim
(\code{radiation\_isrf.c:1362-1369}). Let $\boldsymbol{F}$ be the
\emph{pre-update} flux, snapshotted per operator before the moment loop
overwrites it (\code{radiation\_isrf.c:1278}, \code{:1301-1303}), and let
$\code{reduced\_flux} = \code{isrf\_c\_hyp\_consistent\_variable\_c}$, i.e.\
true for schemes 3 and 4:

\begin{align}
  &\epsilon_R \le 0 \ \Rightarrow\ s = 1; \qquad c_{\rm hyp} \le 0 \ \Rightarrow\ s = 1,
  \label{eq:isrfd-gate-early} \\
  w &= \kappa + \frac{\lambda\, H}{c} \quad \text{(the \textbf{true} $c$, not $c_{\rm hyp}$)},
  \qquad w \le 0 \ \Rightarrow\ s = 1,
  \label{eq:isrfd-gate-w} \\
  c_\nabla &=
    \begin{cases}
      1, & \code{reduced\_flux} \ \text{(schemes 3, 4)},\\
      c_{\rm hyp}, & \text{otherwise},
    \end{cases}
  \label{eq:isrfd-gate-cnabla} \\
  N_R &= \big| w\boldsymbol{F} + c_\nabla \boldsymbol{G} \big|,
  \qquad
  D_R = w\,|\boldsymbol{F}| + c_\nabla\,|\boldsymbol{G}|,
  \qquad D_R \le 0 \ \Rightarrow\ s = 0,
  \label{eq:isrfd-gate-num-den} \\
  R &= \frac{N_R}{D_R},
  \qquad
  s = \min\!\left(1, \left(\frac{R}{\epsilon_R}\right)^{\!2}\right),
  \label{eq:isrfd-gate-s}
\end{align}

at \code{radiation\_isrf.c:1135-1136}, \code{:1144}, \code{:1150},
\code{:1159}, \code{:1162-1165}, \code{:1170}, \code{:1178} and
\code{:1184-1186}, with the scheme flag passed in from the call site at
\code{radiation\_isrf.c:1364-1366}.

Equation~(\ref{eq:isrfd-gate-cnabla}) is the gradient's weight at the discrete
fixed point of the unlimited recurrence
$\boldsymbol{F}^{n+1} = \mathcal{E}\boldsymbol{F} - \mathcal{C}\boldsymbol{G}$.
With $1 - \mathcal{E} = \mathcal{A}\varphi$ and
$\mathcal{A} = c_{\rm hyp} w\, dt$, the two scheme families have different fixed
points, and that is why $c_\nabla$ must switch:

\begin{align}
  \text{schemes 0, 1, 2} \ (\mathcal{C} = c_{\rm hyp}^2 dt\,\varphi): &\quad
    \boldsymbol{F} = -\frac{c_{\rm hyp}}{w}\boldsymbol{G}
    \ \Longleftrightarrow\ w\boldsymbol{F} + c_{\rm hyp}\boldsymbol{G} = 0,
  \label{eq:isrfd-gate-fp-true} \\
  \text{schemes 3, 4} \ (\mathcal{C} = c_{\rm hyp} dt\,\varphi,\ \text{state } \tilde{\boldsymbol{F}}): &\quad
    \tilde{\boldsymbol{F}} = -\frac{1}{w}\boldsymbol{G}
    \ \Longleftrightarrow\ w\tilde{\boldsymbol{F}} + \boldsymbol{G} = 0,
  \label{eq:isrfd-gate-fp-reduced}
\end{align}

per the function's own doxygen (\code{radiation\_isrf.c:1117-1125}). $N_R$, $D_R$,
$R$ and $R^2$ are all formed in \textbf{double}, because the squared denominator
underflows float32 under this build's folding of the ratio and its square into one
division (\code{radiation\_isrf.c:1152-1155}, \code{:1180-1183});
$D_R \le 0$ is branched rather than epsilon-guarded for the same reason
(\code{radiation\_isrf.c:1172-1177}).

One historical note, because it affects measured numbers: an earlier form of
this gate used $c_\nabla = c_{\rm hyp}$ unconditionally, which left one extra
power of $c_{\rm hyp}$ on the gradient term under schemes 3 and 4 and therefore
made the gate inert in the default configuration, scheme~4 being the default
(\code{feedback\_properties.h:1100-1102}), so any floor measurement taken before
the fix is affected. Equations~(\ref{eq:isrfd-gate-cnabla}) to
(\ref{eq:isrfd-gate-fp-reduced}) are the current, corrected form.

\subsubsection{The $c_{\rm hyp}$ schemes 0 to 4}
\label{sec:isrf-eqdisc-chyp}

The enum is \code{isrf\_c\_hyp\_scheme}
(\code{feedback\_properties.h:86-127}). It spans two independent axes: the speed
formula, and whether the pairwise operators use the consistent-variable-$c$ change
of variable. The selector \code{isrf\_c\_hyp\_consistent\_variable\_c} is a
process global, set once (\code{feedback\_properties.h:1120-1122}, defined
\code{radiation\_isrf.c:53}), because the pairwise hooks share the fixed hydro
signature \code{(r2, dx, hi, hj, pi, pj, a, H)} with no \code{engine} pointer
(\code{src/feedback/GEAR/radiation\_isrf.h:73-88}).

With $h_{\rm phys} = a\, h$ (note: \emph{not} $\gamma_K h$),
$dt_i = \code{dt\_prev}$, and $dt_{\max}(i)$ the physical step at
\code{max\_ngb\_time\_bin}:

\begin{center}
\begin{tabular}{p{3.6cm}p{5.3cm}p{1.6cm}p{3.4cm}}
\hline
scheme & speed formula & operators & code \\
\hline
0 \code{shipped} & $c_{\rm hyp} = \min(C_{\rm hyp} h_{\rm phys}/dt_i,\, c)$ & shipped & \code{radiation\_isrf.c:288-289} \\
1 \code{kernel\_local} & $c_{\rm hyp} = \min(C_{\rm hyp} h_{\rm phys}/dt_{\max}(i),\, c)$ & shipped & \code{radiation\_isrf.c:530-531} \\
2 \code{fixed\_fraction} & $c_{\rm hyp} = f_c\, c$ & shipped & \code{radiation\_isrf.c:277-278} \\
3 \code{consistent\_variable\_c} & scheme 0's formula & variable-$c$ & \code{radiation\_isrf.c:288-289} \\
4 \code{kernel\_local\_plus\_variable\_c} & scheme 1's formula & variable-$c$ & \code{radiation\_isrf.c:530-531} \\
\hline
\end{tabular}
\end{center}

Scheme 4 is the \textbf{default} (\code{feedback\_properties.h:1100-1102}), and
schemes 3 and 4 additionally set
\code{isrf\_c\_hyp\_consistent\_variable\_c} (\code{feedback\_properties.h:1120}).
Here $f_c = \code{ISRF\_c\_hyp\_fixed\_fraction\_of\_c}$ and
$C_{\rm hyp} = \code{ISRF\_c\_hyp\_margin}$, default $0.5$
(\code{feedback\_properties.h:1169-1170}).

Schemes 0, 2 and 3 decide $c_{\rm hyp}$ at the drift
(\code{radiation\_isrf.c:266-300}); schemes 1 and 4 defer it to the density ghost,
once the $h$-iteration has converged (\code{radiation\_isrf.c:493-543}).
\code{max\_ngb\_time\_bin} is accumulated unconditionally, one byte compare per
pair (\code{radiation\_propagation\_iact.h:430-431}, \code{:486}), and reset to
the particle's own bin each $h$-iteration (\code{radiation\_isrf.c:438}).

A degenerate $dt \le 0$ returns $c$ \emph{directly} rather than relying on
$\min(+\infty, c)$, which \code{-ffast-math} does not guarantee
(\code{radiation\_isrf.c:279-286}, \code{:522-528}). The debug pin
\code{ISRF\_c\_hyp\_pin\_for\_debugging} is applied \emph{after} the light-speed
clamp and is deliberately not itself clamped, so a superluminal pin is possible
on purpose (\code{radiation\_isrf.c:296-297}, \code{:537-538}).
$dt_{\max}(i)$ reuses \code{dt\_prev} in the same-bin case rather than
recomputing, for bit-identity with scheme 0 there
(\code{radiation\_isrf.c:504-505}). At parse time
\code{feedback\_props\_check\_c\_hyp\_scheme} forces $f_c > 0$ and scheme 2 to
agree (\code{feedback\_properties.h:633-643}), and scheme 2 and the pin are
mutually exclusive (\code{feedback\_properties.h:1158-1164}).

\subsubsection{The timestep each scheme imposes}
\label{sec:isrf-eqdisc-timestep}

Only scheme 2 imposes a radiation timestep.
\code{radiation\_isrf\_part\_timestep} (\code{radiation\_isrf.c:381-417}):

\begin{equation}
  dt_{\rm rad} = \frac{C_{\rm hyp}\, h_{\rm phys}}{f_c\, c}.
  \label{eq:isrfd-dt-rad}
\end{equation}

at \code{radiation\_isrf.c:415-416}. It returns \code{FLT\_MAX}, i.e.\ no
constraint, if $f_c \le 0$ (\code{radiation\_isrf.c:389}), if
\code{ISRF\_propagation} is off (\code{:390}), if the debug off-switch is set
(\code{:391-392}), if the particle is not near-field (\code{:400-405}), or if
$h_{\rm phys} \le 0$ (\code{:413}, an explicit branch rather than a clamp,
because a clamp does not shield a NaN under \code{-ffast-math}). The near-field
set (\code{radiation\_isrf.c:400-404}) is: \code{is\_illuminated\_ISRF}, or
$u_{\rm PE} \ne 0$, or $u_{\rm LW} \ne 0$, or either operator's
\code{ngb\_mean\_abs\_u\_V} $>0$. \code{ISRF\_MOMENT\_LW\_PHOTON} is deliberately
omitted, because first init forces it nonzero if and only if $u_{\rm LW}$ is
(\code{radiation\_isrf.c:395-399}).
\code{feedback\_compute\_part\_timestep} additionally aborts the run if the
scheme-2 bound falls below \code{TimeIntegration:dt\_min}, naming the parameter
(\code{feedback.c:194-212}).

Schemes 0, 1, 3 and 4 impose \textbf{no} radiation timestep at all: their
$c_{\rm hyp}$ is already derived from a timestep
(\code{radiation\_isrf.c:345-356}). The coverage is not the same for all four.
Schemes 1 and 4 build $c_{\rm hyp}$ from $dt_{\max}(i)$, so the receiver bound
$c_i\, dt_j \le C_{\rm hyp} h_i$ holds for every neighbour the density loop
reached, i.e.\ every $j$ inside $\mathcal{H}_i$
(\code{feedback\_struct.h:450-451}). Schemes 0 and 3 build it from $dt_i$ alone,
so the bound holds only for a neighbour on this particle's own bin and is
violated by one on a coarser bin; that gap is the stated reason scheme 1 exists.
Neither case covers the $\mathcal{H}_i \le r < \mathcal{H}_j$ pair class
(Section~\ref{sec:isrf-eqdisc-limits}, item~\ref{it:isrfd-uncovered}).

Absorption imposes no timestep limit by construction: the exponential relaxation
is exact at any $\kappa\, dt$.

\subsubsection{The dose reservoir and its drawdown}
\label{sec:isrfd-reservoir}

Drawdown happens at the drift, for active particles only
(\code{radiation\_isrf.c:318-340}). With
$\code{ti\_begin} = \code{get\_integer\_time\_begin}(\code{ti\_current}, \code{time\_bin})$:

\begin{align}
  t_{\rm rem} &=
    \begin{cases}
      0, & \code{ISRF\_reservoir\_end\_ti} \le \code{ti\_begin},\\
      \code{cosmology\_get\_delta\_time}(\code{ti\_begin}, \code{..\_end\_ti}), & \text{with cosmology},\\
      (\code{..\_end\_ti} - \code{ti\_begin})\,\code{time\_base}, & \text{otherwise},
    \end{cases}
  \label{eq:isrfd-trem} \\
  f_{\rm draw} &=
    \begin{cases}
      1, & t_{\rm rem} \le dt_{\rm phys},\\
      dt_{\rm phys}/t_{\rm rem}, & \text{otherwise},
    \end{cases}
  \label{eq:isrfd-fdraw} \\
  S_m &= \frac{f_{\rm draw}\, \code{u\_dose\_reservoir}_m}{dt_{\rm phys}},
  \qquad
  \code{u\_dose\_reservoir}_m \mathrel{-}= f_{\rm draw}\, \code{u\_dose\_reservoir}_m,
  \label{eq:isrfd-drawdown}
\end{align}

at \code{radiation\_isrf.c:320-331} and \code{:334-335}. Passing
\code{ti\_current} itself instead of \code{ti\_begin} would over-drain and empty
the reservoir one sub-step early (\code{radiation\_isrf.c:243-248}), and
$dt_{\rm phys} > 0$ is required so the refund below is exact
(\code{radiation\_isrf.c:303-316}). A particle whose $h$-iteration gives up with
no neighbours is refunded exactly the amount subtracted:
$\code{u\_dose\_reservoir} \mathrel{+}= S_m\, \code{dt\_prev}$ with
$S_m \leftarrow 0$ (\code{radiation\_isrf.c:975-976}).

\subsection{Injection}
\label{sec:isrf-eqdisc-injection}

\code{radiation\_iact\_nonsym\_feedback\_apply}
(\code{radiation\_iact.h:202-355}). The gate is
$L_{\rm band}(\mathrm{PE}) \ne 0$ or $L_{\rm band}(\mathrm{LW}) \ne 0$
(\code{radiation\_iact.h:280-281}).

\subsubsection{Kernel weighting and normalisation}

\begin{align}
  r &= \sqrt{\max\big(r^2,\ 10^{-6} h_i^2\big)},
  \label{eq:isrfd-inj-r} \\
  w_i &= W(r/h_i)\, h_i^{-d},
  \label{eq:isrfd-inj-w} \\
  \omega_{ij} &= \frac{m_j\, w_i}{\code{enrichment\_weight}_i},
  \label{eq:isrfd-inj-weight}
\end{align}

at \code{radiation\_iact.h:214-215}, \code{:218-222} and \code{:224-227}, with
$r$ and $h_i$ \textbf{comoving} and $d$ the dimension count. Here
\code{enrichment\_weight} is $\sum_j m_j W(r/h_i)$ accumulated in the star
density loop (\code{feedback\_iact.h:71}) and multiplied by $h_i^{-d}$ in
\code{feedback\_prepare\_feedback} (\code{feedback.c:366}), i.e.\ it \emph{is} the
SPH density estimate at the star; a value of 0 gives
$\omega_{ij} = 0$ rather than a division
(\code{radiation\_iact.h:224-226}). The injected energy per pair, per moment, is

\begin{equation}
  \Delta U_m = \Delta t_\star\, \omega_{ij}\, L_{\rm band}(m)\,
    \mathcal{X}_{o(m)},
  \label{eq:isrfd-inj-energy}
\end{equation}

at \code{radiation\_iact.h:296-300}, with $\mathcal{X}$ the extinction factor of
Section~\ref{sec:isrf-eqdisc-extinction}. $\Delta t_\star$ is the star's own
cached feedback step in proper time (\code{radiation\_iact.h:236}, cached
\code{:124}), asserted against a stale step boundary under
\code{SWIFT\_DEBUG\_CHECKS} (\code{radiation\_iact.h:237-244}).

The deposit depends on the mode:

\begin{itemize}
  \item \code{ISRF\_propagation} \textbf{on}
    (\code{radiation\_iact.h:309-315}): pure accumulation into the reservoir,
    $\code{u\_dose\_reservoir}_m \mathrel{+}= \Delta U_m / m_j$ in float, plus
    $\code{ISRF\_reservoir\_end\_ti} = \max(\text{old},\
    \code{ti\_current} + \code{ti\_step\_star})$. There is no reset and no
    first-touch logic, so any number of stars on any time bins superpose.
  \item \code{ISRF\_propagation} \textbf{off}
    (\code{radiation\_iact.h:322-333}): $u$ is an \emph{instantaneous}
    strength. On the first touch of the step every band's $u$ is zeroed
    (\code{radiation\_iact.h:322-326}), then
    $u_m \mathrel{+}= \Delta U_m/m_j$ in double, so
    \begin{equation}
      u_m = \frac{1}{m_j}\sum_{\star} \Delta t_\star\, \omega_{\star j}\,
        L_{\rm band}(m)\, \mathcal{X}_{o(m)},
      \label{eq:isrfd-inj-instantaneous}
    \end{equation}
    at \code{radiation\_iact.h:328-333}, which scales linearly with the
    illuminating star's own timestep, and so does $G_0$.
\end{itemize}

The illumination window is renewed on every touch,
$\code{ISRF\_illumination\_end\_ti} = \code{ti\_current} + 2\,\code{ti\_step\_star}$
(\code{radiation\_iact.h:342-343},
$\code{RADIATION\_ISRF\_TAG\_LIFETIME\_INTERVALS} = 2$,
\code{radiation.h:61}), and \code{timestep\_sync\_part} is called on the first
touch only (\code{radiation\_iact.h:350-353}).

\subsubsection{Dust extinction factor}
\label{sec:isrf-eqdisc-extinction}

\begin{align}
  \frac{\mathcal{D}(Z)}{\mathcal{D}_{\rm MW}} &=
    \frac{\max(Z, 0)}{0.01295}\cdot\frac{f_{\rm gr}}{0.009387},
  \label{eq:isrfd-dust-ratio} \\
  \kappa_{\rm eff} &= \frac{\sigma_d(b)\,\mathcal{D}(Z)/\mathcal{D}_{\rm MW}}
    {\mu_H\, m_p\, \code{units\_cgs}(\mathrm{AREA})},
  \label{eq:isrfd-kappa-eff} \\
  \mathcal{X} &= \exp\big(-\kappa_{\rm eff}\, \Sigma_{\rm gas}\big),
  \label{eq:isrfd-extinction} \\
  \Sigma_{\rm gas} &= \ell_{\rm ext}\, \rho\, a^{-2},
  \label{eq:isrfd-column}
\end{align}

at \code{radiation\_isrf.c:1500-1503}, \code{:1526-1529}, \code{:1552} and
\code{:1480-1481}, with the $a^{-2}$ applied at
\code{radiation\_isrf.c:1604-1605}. In Equation~(\ref{eq:isrfd-column}) the
$\rho$ is the receiver's \textbf{comoving} \code{p->rho} and $\ell_{\rm ext}$ is
comoving, so the $a^{-2}$ makes $\Sigma_{\rm gas}$ a \textbf{physical} column, in
mass per area.

$\kappa_{\rm eff}$ is a \textbf{mass opacity} (area per mass, internal units) and
must be kept distinct from the per-particle $\kappa$ of
Section~\ref{sec:isrf-eqdisc-continuum}, which is a \emph{linear} rate
(1/length) equal to $\kappa_{\rm eff}\,\rho_{\rm phys}$
(\code{radiation\_isrf.c:1577-1579}). The two are computed independently and the
code states they are not interchangeable (\code{feedback\_struct.h:308-310}).
$f_{\rm gr}$ is the \emph{resolved}
\code{chemistry\_data.local\_dust\_to\_gas\_ratio}, never the $-1$ sentinel
(\code{radiation\_isrf.c:1609-1610}, \code{:225-226}). The factor is
band-specific and computed per operator
(\code{radiation\_isrf.c:1612-1614}).

\subsubsection{Extinction path: three mechanisms}

\code{radiation\_get\_comoving\_extinction\_path}
(\code{radiation\_isrf.c:1416-1463}), with
$h_{\rm gas} = \gamma_K\, h$ (\code{radiation\_isrf.c:1425}). The default
mechanism is \code{pair\_separation}
(\code{feedback\_properties.h:1047-1048}).

\begin{align}
  \text{\code{constant\_kernel\_path}:}\quad
    \ell_{\rm ext} &= R_{\rm ext}\, \gamma_K\, h_j,
  \label{eq:isrfd-path-constant} \\
  \text{\code{pair\_separation} (default):}\quad
    \ell_{\rm ext} &= r,
  \label{eq:isrfd-path-pair} \\
  \text{\code{temperature\_capped\_jeans}:}\quad
    \ell_{\rm ext} &= \min\!\left(\lambda_{\rm J}\, a^{-1},\ \gamma_K h_j\right),
  \label{eq:isrfd-path-jeans}
\end{align}

at \code{radiation\_isrf.c:1431}, \code{:1433-1434} and \code{:1457}, with
$R_{\rm ext}$ the parameter \code{ISRF\_extinction\_path\_in\_kernel\_radii}, default 1.0
(\code{feedback\_properties.h:41}, \code{:1053-1055}), and the Jeans length built
from a temperature-capped sound speed:

\begin{equation}
  \theta = \min\!\left(1, \frac{T_{\rm cap}}{T}\right),
  \qquad
  c_s^2\big|_{\rm capped} = c_s^2\, \theta,
  \qquad
  \lambda_{\rm J} = \sqrt{\frac{\pi\, c_s^2\big|_{\rm capped}}{G\, \rho_{\rm phys}}},
  \label{eq:isrfd-jeans}
\end{equation}

at \code{radiation\_isrf.c:1450-1452} and \code{:1455-1456}, with
$T_{\rm cap} = 40$\,K by default (\code{feedback\_properties.h:43}).

Guards: $T \le 0$ or $\rho_{\rm phys} \le 0$ returns $h_{\rm gas}$ outright,
rather than clamping a division, because a guarded denominator is not safe under
\code{-ffast-math} (\code{radiation\_isrf.c:1441-1444}). An unrecognised
mechanism string is fatal (\code{feedback\_properties.h:1075-1087}), and an
unknown enum value reaching the switch is fatal
(\code{radiation\_isrf.c:1461-1462}). Each mechanism's magnitude key is parsed
only by the mechanism that reads it, so \code{used\_parameters.yml} records a
number only where it was used (\code{feedback\_properties.h:1043-1044}). The
comment calls $R_{\rm ext} = 5/12$ the exact value in the uniform optically thin
limit (\code{feedback\_properties.h:146-149}).

\subsubsection{The exact identity at zero metallicity}
\label{sec:isrfd-zero-Z}

At $Z = 0$, Equation~(\ref{eq:isrfd-dust-ratio}) gives
$\mathcal{D}(Z)/\mathcal{D}_{\rm MW} = 0$
(\code{radiation\_isrf.c:1500-1501}), hence $\kappa_{\rm eff} = 0$ and
$\mathcal{X} = e^{0} = 1$ exactly for both bands. Summing the deposited
\emph{energy} over the star's neighbours then gives

\begin{equation}
  \sum_j m_j\,\frac{\Delta U_m}{m_j}
  = \Delta t_\star\, L_{\rm band}(m) \sum_j \frac{m_j w_i}{\rho_\star}
  = \Delta t_\star\, L_{\rm band}(m),
  \label{eq:isrfd-inj-identity}
\end{equation}

because $\rho_\star = \sum_j m_j w_i$ is the same kernel sum
(\code{feedback\_iact.h:71} with \code{feedback.c:366}). So the injection
conserves the star's band energy over the step exactly, with three documented
exceptions:

\begin{itemize}
  \item the apply loop floors $r$ at $10^{-6}h_i^2$
    (\code{radiation\_iact.h:214-215}) while the density loop that built
    \code{enrichment\_weight} does not (\code{feedback\_iact.h:58}), so a
    coincident pair breaks the identity;
  \item the reservoir accumulator is \textbf{float}
    (\code{radiation\_iact.h:311}), unlike $u$ itself;
  \item the identity holds for the \emph{deposited dose}. What reaches $u$ is
    $(c_{\rm hyp}/c)\times\text{dose}$ plus the relaxation factors
    (\code{radiation\_isrf.c:936}): injection deposits the raw, unrescaled dose
    and the rescale is applied exclusively in the end-force update
    (\code{radiation\_isrf.c:717-720}).
\end{itemize}

Note also that \code{radiation\_relaxation\_phi\_factor}'s doxygen names an
``injection-site equivalent, $c_{\rm hyp}\kappa\Delta t_\star$''
(\code{radiation\_isrf.c:558-559}). No such caller exists: the only callers are
the two ghosts inside \code{radiation\_isrf.c}. The injection applies no
$\varphi$.

\subsection{The band-edge transfer, LW into PE}
\label{sec:isrf-eqdisc-transfer}

Computed before the moment loop, from LW's own depth and frozen inputs, so the
result cannot depend on moment ordering
(\code{radiation\_isrf.c:834-886}). The whole block is gated on
$H_{\rm dil} \ne 0$ (\code{radiation\_isrf.c:835}):

\begin{align}
  \mathcal{A}_{\rm LW} &= \big(c_{\rm hyp}\kappa_{\rm LW}
    + \lambda_E(\mathrm{LW})\, H_{\rm dil}\big)\, dt,
  \label{eq:isrfd-transfer-depth} \\
  1 - \mathcal{E}_{\rm LW} &= -\code{expm1}(-\mathcal{A}_{\rm LW}),
  \qquad
  \varphi_{\rm LW} = \varphi(\mathcal{A}_{\rm LW}) \ \text{[double]},
  \label{eq:isrfd-transfer-decay} \\
  \mathcal{L}_{\rm LW} &=
    \Big(u_{{\rm prev},{\rm LW}} + dt\,\varphi_{\rm LW}\,
      \code{dissipation\_u}_{\rm LW}\Big)\big(1 - \mathcal{E}_{\rm LW}\big)
    \nonumber\\
    &\qquad + \Big(\mathcal{R}\, S_{\rm LW} - X_{\rm LW}\Big)\,
      dt\,\big(1 - \varphi_{\rm LW}\big),
  \label{eq:isrfd-transfer-absorbed} \\
  f_{\rm edge} &= \frac{\big(\lambda_E(\mathrm{LW}) - 1\big)\, H_{\rm dil}\, dt}
    {\mathcal{A}_{\rm LW}}
    \qquad \text{if } \mathcal{A}_{\rm LW} > 0,
  \label{eq:isrfd-transfer-fedge} \\
  \mathcal{T} &= f_{\rm edge}\, \mathcal{L}_{\rm LW},
  \label{eq:isrfd-transfer}
\end{align}

at \code{radiation\_isrf.c:840-842}, \code{:849-850}, \code{:857-862},
\code{:870-871} and \code{:884}. The result is added to PE by
Equation~(\ref{eq:isrfd-u-transfer-add}) at \code{radiation\_isrf.c:944}, with no
$\varphi$ and no PE absorption.

\paragraph{What it conserves.} $\mathcal{L}_{\rm LW}$ is the total amount that
LW's relaxation depth removed this step, summed over all competing sinks in that
depth, i.e.\ dust absorption plus the band-edge loss. $f_{\rm edge}$ is exactly
the band-edge loss's share of $\mathcal{A}_{\rm LW}$, since
$\mathcal{A}_{\rm LW} = (c_{\rm hyp}\kappa + \lambda_E(\mathrm{LW})H_{\rm dil})dt$
and the band-edge part of it is
$(\lambda_E(\mathrm{LW}) - 1) H_{\rm dil} dt$. The transfer is therefore an
\textbf{exact proportional split} of one combined exponential sink, and LW's loss
at its $11.2$\,eV lower edge equals PE's gain at its $11.2$\,eV upper edge
exactly, because the two bands are contiguous there
(\code{radiation\_isrf.c:723-736}). $f_{\rm edge}$ is derived from the
\emph{same} $\mathcal{A}_{\rm LW}$ the moment loop uses, not from a separately
written algebraically-equal expression, because under \code{-ffast-math} two such
expressions can disagree at extreme $\lambda$ and let the loss apply while the
gain silently does not (\code{radiation\_isrf.c:816-821}). Grey LW,
$\lambda_E(\mathrm{LW}) = 1$, gives $\mathcal{T} = 0$ exactly.

\paragraph{Deliberate asymmetries and limits.}

\begin{itemize}
  \item The transfer is added raw to PE's already-relaxed $u$, so it is not
    subject to PE's own absorption during its arrival step
    (\code{radiation\_isrf.c:939-944}).
  \item $\mathcal{T}$ is kept \textbf{signed}, not clamped at zero: clamping
    would break the exact $11.2$\,eV cancellation and the per-band ledger. A
    negative transfer (LW undershot) reduces PE's $u$ on the same step and can
    drive it negative, which then zeroes PE's flux on the next step's limiter
    (\code{radiation\_isrf.c:872-883}).
  \item \code{ISRF\_MOMENT\_LW\_PHOTON} loses at its own lower edge, with
    $\lambda_N(\mathrm{LW})$, but has no photon-number PE moment to transfer
    into: that loss is uncompensated (\code{radiation\_isrf.c:730-735}).
  \item The PE band's own $6$\,eV lower-edge loss is kept, not compensated, by
    design (\code{radiation\_isrf.c:735-736}).
  \item $\mathcal{A}_{\rm LW} > 0$ is tested exactly before the division in
    Equation~(\ref{eq:isrfd-transfer-fedge}), to avoid $0/0$ for a metal-free
    particle with $H_{\rm dil} \ne 0$ (\code{radiation\_isrf.c:863-869}).
  \item The $H_{\rm dil} \ne 0$ gate is the physically correct statement and the
    cheaper non-cosmological path, but the comment explicitly denies that it is
    a bit-identity guarantee across two builds under
    \code{-flto -O3 -ffast-math} (\code{radiation\_isrf.c:823-833}).
\end{itemize}

\subsection{Couplings out to the gas}
\label{sec:isrf-eqdisc-couplings}

\subsubsection{Habing $G_0$, the sum of both bands}

\code{radiation\_get\_part\_isrf\_habing} (\code{radiation\_gas.c:761-780}):

\begin{align}
  u_\Sigma &= \big[u_{\rm PE}\big]_+ + \big[u_{\rm LW}\big]_+,
  \label{eq:isrfd-usum} \\
  \mathcal{J} &= c\, \rho_{\rm phys}\, u_\Sigma,
  \label{eq:isrfd-habing-flux} \\
  G_0 &= \left[\frac{\mathcal{J}\,\code{units\_cgs}(\text{energy flux per unit surface})}
    {1.6\times10^{-3}}\right]_+,
  \label{eq:isrfd-G0}
\end{align}

at \code{radiation\_gas.c:767-768}, \code{:769}, \code{:770-772} and
\code{:777-779}, where $[\cdot]_+$ is the non-negativity clamp.

$G_0$ is the \textbf{sum of both bands}, because the Habing band is $6$ to
$13.6$\,eV, which is exactly PE plus LW.
\code{ISRF\_MOMENT\_LW\_PHOTON} does not enter. The $c$ in
Equation~(\ref{eq:isrfd-habing-flux}) is the \textbf{true} speed of light, not
$c_{\rm hyp}$, and $\rho$ is physical. Each band is clamped non-negative
\emph{before} the sum (\code{radiation\_get\_band\_u\_nonnegative},
\code{radiation\_gas.c:726-736}), so an undershot band reads as zero
illumination instead of cancelling part of the other band's real signal, and the
summed result passes the same clamp again. The clamp is read-only on the
particle: the stored $u$ keeps its negative value
(\code{radiation\_gas.c:718-720}).

$G_0$ reaches Grackle's per-particle \code{isrf\_habing} only when
\code{chemistry\_data.use\_isrf\_field} is set (\code{cooling.c:932-937}), which
\code{with\_interstellar\_radiation\_field} forces on internally together with
\code{dust\_chemistry} and \code{photoelectric\_heating}
(\code{cooling\_properties.h:123-131}).
\code{cooling\_get\_isrf\_habing\_subgrid} carries \textbf{no} Grackle-mode gate
(\code{cooling\_gear\_subgrid.h:267-273}): it returns 0 only if the ISRF is off.

\subsubsection{H$_2$ Lyman-Werner photodissociation rate}

\code{radiation\_get\_part\_LW\_dissociation\_rate\_internal}
(\code{radiation\_gas.c:816-837}):

\begin{align}
  \mathcal{J}_{\rm LW} &= c\, \rho_{\rm phys}\, \big[u_{\rm LW}\big]_+,
  \label{eq:isrfd-lw-flux} \\
  k_{\rm diss} &= \left[\frac{\dfrac{\sigma_{\rm H_2}}{E_{\rm LW}}\,
    \mathcal{J}_{\rm LW}\,\code{units\_cgs}(\text{energy flux per unit surface})}
    {\code{units\_cgs}(\mathrm{INV\_TIME})}\right]_+,
  \label{eq:isrfd-kdiss}
\end{align}

at \code{radiation\_gas.c:821-822}, \code{:823-825}, \code{:827} and
\code{:828-834}, with $\sigma_{\rm H_2}/E_{\rm LW}$ the single constant of
Equation~(\ref{eq:isrfd-sigma-over-e}). The cross section and the photon energy
never appear separately, and \textbf{no} table value enters the rate:
\code{radiation\_lw\_photon\_energy\_cgs} is a reported diagnostic only
(\code{radiation\_gas.c:792-797},
\code{src/feedback/GEAR/radiation\_isrf.h:43-52}).
\code{ISRF\_MOMENT\_LW\_PHOTON} does not enter either.

\paragraph{The \texttt{COOLING\_GRACKLE\_MODE > 1} gate.}
\code{cooling\_get\_LW\_dissociation\_rate\_subgrid}
(\code{cooling\_gear\_subgrid.h:288-299}) returns
Equation~(\ref{eq:isrfd-kdiss}) only under
\code{\#if COOLING\_GRACKLE\_MODE > 1}, and $0$ otherwise. So at
\code{grackle\_0} and \code{grackle\_1} the rate is \textbf{identically zero},
with no start-up warning. The LW band is still transported and still written to
snapshots, and its \emph{energy} still feeds $G_0$ through
Equation~(\ref{eq:isrfd-usum}), so what is inert below mode 2 is only the band's
own dedicated dissociation effect.

Two further facts at the Grackle interface. The per-particle rate takes priority
over the global YAML scalar only when it is nonzero
(\code{cooling.c:905-916}), so at mode 2 or above a particle with
$k_{\rm diss} = 0$ (unilluminated, or clamped) receives the global
\code{GrackleCooling:RT\_H2\_dissociation\_rate\_cgs} instead of zero; setting
both is only a warning (\code{cooling\_io.h:325-333}). And the per-particle rate
reaches Grackle \emph{unshielded} unless
\code{GrackleCooling:H2\_self\_shielding} selects a column-length mode; the code
warns that at $N_{\rm H_2} = 10^{20}\,{\rm cm^{-2}}$ the
Wolcott-Green \& Haiman (2019) factor is of order $10^{-5}$, so H$_2$ cannot survive
an illuminated cloud without it (\code{cooling\_io.h:417-430}).

\subsection{Approximations and known limitations}
\label{sec:isrf-eqdisc-limits}

\subsubsection{Physical approximations}

\begin{enumerate}
  \item \textbf{Reduced speed of light.} $c_{\rm hyp} \le c$ replaces $c$ in
    every transport and absorption rate, and the source is rescaled by
    $c_{\rm hyp}/c$ (\code{radiation\_isrf.c:936}). The steady state is
    preserved by the rescale, but every \emph{transient} is slowed by
    $c_{\rm hyp}/c$, including the pure-expansion transient at $\kappa = 0$,
    which decays at $(c_{\rm hyp}/c)H$ instead of $H$
    (\code{radiation\_isrf.c:712-717}). Where the light-speed clamp is not
    active, $c_{\rm hyp}\, dt = C_{\rm hyp}\, h_{\rm phys}$ exactly, so the
    kernel light-crossing time is pinned to a fixed fraction of a timestep
    regardless of the true light travel time.
  \item \textbf{Grey within each band.} One dust cross section per band
    (\code{radiation.h:105-106}) and one spectrum-averaged
    $\sigma_{\rm H_2}/E_{\rm LW}$ (\code{radiation.h:147}), both compile-time
    constants rather than per-run parameters (\code{radiation.h:156-166}).
  \item \textbf{Run-wide scalar band-edge weights}, evaluated once over the
    whole IMF at one reference metallicity $Z = 0.014$
    (\code{radiation.c:122-127}). The comment states that the bias of the
    shipped scalar against a population average is \emph{not sized}
    (\code{radiation.h:230-236}), and that an aged population's
    $\Lambda_{\rm LW}$ reaches about 18.3, roughly three times the shipped
    young-population value.
  \item $\mu_H = 1.4$ and $Z_\odot = 0.01295$ are calibration partners of the
    grain model, not the run's own composition (\code{radiation.h:94-100},
    \code{:108-114}).
  \item \textbf{Receiver-side extinction column with no resolved density
    gradient.} The column is the receiver's own local density times a path
    length (\code{radiation\_isrf.c:1480-1481}), and every path mechanism
    assumes the intervening medium sits at the receiver's density
    (\code{feedback\_properties.h:155-157}).
  \item The transported \code{LW\_PHOTON} moment is energy-equivalent at a fixed
    reference photon energy (\code{feedback\_struct.h:742-744}), so a spectrum's
    own $E_{\rm LW}$ cannot change the dissociation rate
    (\code{radiation\_gas.c:794-797}).
  \item \textbf{No $v/c$ terms anywhere.} $u$ and $\boldsymbol{F}$ are
    mass-specific quantities on Lagrangian particles, so the
    comoving-frame/lab-frame distinction is not carried. This is an observation
    from reading the equations, not a code comment.
\end{enumerate}

\subsubsection{Clamps, floors and epsilon guards}

Each item below is a place where the scheme departs from its own ideal form.

\begin{enumerate}
  \setcounter{enumi}{7}
  \item $\exp(-\kappa_{\rm eff}\Sigma)$ has no column cap, and $\max(Z,0)$
    floors a negative metallicity (\code{radiation\_isrf.c:1500-1501}).
  \item The M1 flux limiter, Equation~(\ref{eq:isrfd-limiter}): its
    $u \le 0$ branch zeroes a band's flux entirely.
  \item The $\varphi$ Taylor branch, Equation~(\ref{eq:isrfd-phi}).
  \item \code{rho\_prev} floored to $1.0$ when the comoving density is
    non-positive (\code{radiation\_isrf.c:213}) and seeded to $1.0$ at first
    init (\code{radiation\_isrf.c:158}), to stop $1/\rho$ becoming $+\infty$.
  \item The $10^{-6}h_i^2$ floor on the injection pair separation
    (\code{radiation\_iact.h:214-215}, \code{:75-76}).
  \item The light-speed clamp and the $dt \le 0$ branch of
    Section~\ref{sec:isrf-eqdisc-chyp}.
  \item In the dissipation: $x = \min(\epsilon/\epsilon_1, 1)$,
    Equation~(\ref{eq:isrfd-trig-aim}); the quartic roll-off,
    Equation~(\ref{eq:isrfd-floor}); and the gate's
    $\min(1, (R/\epsilon_R)^2)$ with its four degenerate early returns,
    Equations~(\ref{eq:isrfd-gate-early}) to (\ref{eq:isrfd-gate-s}).
  \item Non-negativity clamps at the Grackle interface, per band and again on
    the summed result (\code{radiation\_gas.c:665-691}, \code{:726-736}). The
    code calls these ``the last line of defence \ldots not a fix for the
    underlying undershoot'' (\code{radiation\_gas.c:656-657}) and warns,
    throttled to the first 10 events then every 10000
    (\code{radiation\_gas.c:645-646}, \code{:680-688}).
  \item \code{radiation\_get\_band\_u\_nonnegative} is read-only on the
    particle, so the stored $u$ and the \code{u\_min\_since\_snapshot}
    diagnostic keep the negative value (\code{radiation\_gas.c:718-720}): the
    undershoot is hidden from Grackle but visible in the snapshot.
\end{enumerate}

\subsubsection{Numerical and scheme limitations the code itself flags}

\begin{enumerate}
  \setcounter{enumi}{16}
  \item \label{it:isrfd-uncovered} \textbf{The uncovered pair class
    $\mathcal{H}_i \le r < \mathcal{H}_j$} (\code{feedback\_struct.h:452-477}).
    The force loop fires a side whenever either kernel reaches, and the shared
    $\Phi_{ij}$ of Equation~(\ref{eq:isrfd-div-phi}) carries a nonzero
    $W'_j$ into particle $i$ even where $W'_i = 0$; but the
    density loop never saw $j$, so $dt_j$ never entered $dt_{\max}(i)$. The
    overshoot factor is $dt_j/dt_{\max}(i)$, unbounded while $j$ stays
    inactive; once $j$ recomputes its own step the
    $\code{min\_ngb\_bin} + 2$ cap bounds it at 4. Closing the gap needs
    $dt_{\max}(i)$ built over the force loop's neighbour set.
  \item \textbf{Stability bound with hard-coded kernel constants}
    (\code{feedback\_properties.h:1254-1275}):
    \begin{equation}
      6.2\,\alpha\, C_{\rm hyp} + 0.70\, C_{\rm hyp}^2 \le 2,
      \label{eq:isrfd-stability}
    \end{equation}
    with $6.2 = 2 I_W$ and $0.70 = \nu_{\max}^2/2$ for Wendland-C2
    ($I_W = 3.10$, $\nu_{\max} = 1.18$). It is \emph{not} re-derived for any
    other kernel. At $\alpha = 0$ it gives the absolute ceiling
    $C_{\rm hyp} \le \sqrt{2/0.70} \approx 1.6903$
    (\code{feedback\_properties.h:1179}); at the shipped
    $\alpha_{\rm floor} = 0.5$ it caps $C_{\rm hyp}$ at about $0.571$
    (\code{feedback\_properties.h:290-294}).
  \item The staggered order is forced: $\boldsymbol{F}$ is advanced first, then
    $u$ from the new flux. The code attributes the order to conservation, not
    to stability (\code{radiation\_isrf.c:693-699}).
  \item The negativity trigger is one step late by construction: it reads
    $u^n$, so an undershoot created by this step's update is corrected by the
    \emph{next} step, and cooling reads it uncorrected for one step through its
    own clamp (\code{radiation\_isrf.c:738-743}).
  \item The dissipation, built from $u^n$, is decayed together with
    $u_{\rm prev}$ rather than added undamped, for the larger stability margin
    (\code{radiation\_isrf.c:697-700}).
  \item $c_{\rm hyp}$ uses $a\, h$, not $\gamma_K h$
    (\code{radiation\_isrf.c:269}, \code{:520}, \code{:407}), while the
    extinction path uses $\gamma_K h$ (\code{radiation\_isrf.c:1425}).
  \item Transport conservation is exact per pair at the \emph{accumulator}
    level only, and only within one time bin and at $\varphi = 1$; see
    Section~\ref{sec:isrf-eqdisc-pairwise}.
  \item The $u$-path and $\boldsymbol{F}$-path relax at different precision; see
    Section~\ref{sec:isrf-eqdisc-F-update}.
  \item MPI and restart: operator and moment state lives on
    \code{struct part} (\code{feedback\_struct.h:415-420}), so it rides the
    normal exchange and dump; the HII owner-computed payload does not
    (\code{feedback\_struct.h:579-598}).
  \item Not independently verified here, and reported as the code's own claims:
    the exact skew-adjointness of the divergence/gradient pair
    (\code{radiation\_propagation\_iact.h:338-341}), and the $6.2$ and $0.70$
    lattice constants of Equation~(\ref{eq:isrfd-stability}). The code names
    its own verification material, \code{theory/GEAR/Radiation/02\_fuv\_isrf.tex}
    and \code{theory/GEAR/Radiation/verify\_isrf\_dissipation.py} (Part~C),
    neither of which lives in the \code{swiftsim} checkout.
\end{enumerate}

\subsubsection{Comment and code disagreements: trust the code}

Each of the following is a place where a doxygen block or inline comment states
something the code does not do. The equations in this section follow the code.

\begin{enumerate}
  \setcounter{enumi}{26}
  \item \code{grad\_u} and the ``$\nabla u$'' wording at
    \code{radiation\_isrf.c:1210} name a gradient of $u$; the accumulator is the
    M1 pressure-tensor divergence, Equation~(\ref{eq:isrfd-grad-apply}).
  \item \code{radiation\_isrf.c:558-559} names an ``injection-site equivalent''
    relaxation depth $c_{\rm hyp}\kappa\Delta t_\star$. No such caller exists:
    the only callers of $\varphi$ are \code{radiation\_isrf.c:850},
    \code{:899} and \code{:1291}. The same doxygen gives
    $\mathcal{A} = c_{\rm hyp}(\kappa + H/c)dt$, omitting $\lambda_m$; the code
    is Equation~(\ref{eq:isrfd-depth-u}).
  \item \code{feedback\_struct.h:277-278} gives
    $\mathcal{A} = (c_{\rm hyp}\kappa + H)\code{dt\_prev}$ for the
    \code{cumulative\_absorbed} ledger: no $c_{\rm hyp}/c$ dilation on $H$, and
    no $\lambda_m$. The code is Equation~(\ref{eq:isrfd-depth-u}).
  \item \code{feedback\_properties.h:384-385} and \code{:403} give
    $C = c_{\rm hyp}^2/(c_{\rm hyp}\kappa + H)$ and
    $w = \kappa + H/c_{\rm hyp}$ for the floor gate. The code uses
    $w = \kappa + \lambda H/c$ with the \textbf{true} $c$
    (Equation~(\ref{eq:isrfd-gate-w}), \code{radiation\_isrf.c:1144}), and $C$
    is $c_{\rm hyp}/w$ per the function's own doxygen
    (\code{radiation\_isrf.c:1071}). The three statements do not agree; the
    code is the last word. The same doxygen's residual, at
    \code{radiation\_isrf.c:1080-1081}, is also still written with
    $c_{\rm hyp}$ on the gradient term unconditionally, i.e.\ it states only
    the true-flux case of Equation~(\ref{eq:isrfd-gate-cnabla}).
  \item \code{radiation\_isrf.c:1210-1216} writes the flux equation as
    $d\boldsymbol{F}/dt = -\boldsymbol{F}/\tau - H\boldsymbol{F}
    - (\mathsf{D}/\tau)\nabla u$ with $\mathsf{D}/\tau = c_{\rm hyp}^2$. In the
    code the spatial term is the M1 pressure-tensor divergence
    (\code{radiation\_propagation\_iact.h:373-388}) and $\mathsf{D}$ is the
    closure tensor, not a diffusion coefficient. That doxygen form is also
    grey-only, with no $\lambda_m$.
  \item Stale output-path references: \code{feedback\_struct.h:146},
    \code{:184} and \code{:202} cite \code{tracers\_io.h} for the snapshot
    fields. The registrations are in fact in
    \code{src/feedback/GEAR\_thermal/feedback\_io.h}.
  \item \code{radiation\_isrf.c:895-896} says of the per-moment band-edge
    weight that ``1 would recover the grey, comoving-edge result this
    replaces''. True for \code{ISRF\_MOMENT\_PE} and \code{ISRF\_MOMENT\_LW};
    \textbf{false} for \code{ISRF\_MOMENT\_LW\_PHOTON}, whose grey value is 0
    by Equation~(\ref{eq:isrfd-lambda-N}). The same wording appears at
    \code{radiation\_isrf.c:726}.
  \item \code{feedback\_struct.h:325} says \code{ngb\_mean\_abs\_u\_V} is
    zeroed ``alongside \code{div\_specific\_flux}''. They are zeroed in
    different places: \code{ngb\_mean\_abs\_u\_V} every $h$-iteration at
    \code{radiation\_isrf.c:440}, \code{div\_specific\_flux} once per step at
    \code{radiation\_isrf.c:1314}.
\end{enumerate}
